\chapter*{Wykaz symboli i skrótów}
\addcontentsline{toc}{chapter}{Wykaz symboli i skrótów}

\begin{description}
  \item[ECAP] Electronic Circuit Analysis Program.
  \item[SPICE] Simulation Program with Integrated Circuit Emphasis; oprogramowanie do modelowania oraz symulacji działania elektronicznych układów analogowych i cyfrowych.
  \item[DC] Direct Current; prąd stały albo analiza stałoprądowa.
  \item[AC] Alternating Current; prąd przemienny albo wymuszenie sinusoidalne.
  \item[ZMPW] Zmodyfikowana Metoda Potencjałów Węzłowych; metoda analizy obwodów elektrycznych oparta na prawach Kirchhoffa i charakterystykach napięciowo-prądowych elementów.
  \item[KCL] Kirchhoff Current Law; pierwsze prawo Kirchhoffa, czyli prądowe prawo Kirchhoffa.
  \item[KVL] Kirchhoff Voltage Law; drugie prawo Kirchhoffa, czyli napięciowe prawo Kirchhoffa.
  \item[BJT] Bipolar Junction Transistor; tranzystor bipolarny.
  \item[UI] User Interface; interfejs użytkownika.
  \item[LAPACK] Linear Algebra Package; biblioteka procedur numerycznych algebry liniowej.
  \item[UTF-8] Unicode Transformation Format 8-bit; kodowanie tekstu używane przez parser.
  \item[UTI] Uniform Type Identifier; systemowy identyfikator typu danych używany w macOS.
  \item[NaN] Not a Number; liczba o niezdefiniowanej lub niereprezentowalnej wartości.
  \item[Inf] Infinity; reprezentacja nieskończoności.
  \item[N-R] Metoda Newtona-Raphsona; algorytm iteracyjny prowadzący do wyznaczenia przybliżonej wartości miejsca zerowego funkcji jednej lub wielu zmiennych.
\end{description}
